\chapter{Costruire e mantenere una coppia di Portali}\label{sec:delivery}

Un incontro fra città richiede che due luoghi mantengano la stessa promessa
nello stesso momento. Quando il collegamento si apre, entrambi devono
avere una persona presente, uno strumento funzionante e qualcuno capace
di accogliere il pubblico. La continuità del servizio è quindi una proprietà
della coppia: un nodo disponibile non basta a garantire l'incontro.
I ruoli e gli accordi qui descritti sono proposte, non incarichi assegnati.

\section{Governare l'incontro da entrambi i lati}
Il responsabile comune del programma concorda obiettivi, calendario,
budget e versione tecnica con un referente per ciascun luogo.
Ogni ospitante risponde di disponibilità, accesso, custodia e convivenza
con le altre attività del proprio sito. Il coordinatore tecnico della
coppia tiene la matrice di compatibilità, le misure e le procedure di
ripristino; un referente locale per nodo segue installazione e guasti.
Il coordinatore educativo concorda le attività con i facilitatori.
La gestione dei dati deve avere responsabilità dichiarate fra i soggetti
coinvolti e procedure definite prima della raccolta.

Il calendario riporta entrambi gli orari locali, i fusi e una data comune
non ambigua; considera cambi d'ora e disponibilità reali delle persone.
La finestra di preparazione deve consentire una prova fra i due facilitatori
prima dell'apertura pubblica. Assenza di un facilitatore, mancata risposta
del partner o anomalia tecnica fanno scattare la procedura concordata:
rinvio, annullamento o passaggio esplicito alla modalità studio locale.
Un'ora di attesa davanti a uno schermo vuoto sarebbe un fallimento di
ospitalità prima ancora che di tecnologia.

Ciascun partecipante può chiedere di interrompere la propria partecipazione;
ciascun facilitatore deve poter chiudere immediatamente il collegamento dal
proprio nodo. La chiusura deve essere segnalata anche all'altro lato, senza
richiedere un suo assenso per fermare audio e immagine locali. Si definiscono
in anticipo condotte ammesse, gestione dei comportamenti inappropriati,
contatto per reclami e procedura per incidenti. Gli aspetti autorizzativi,
assicurativi e di utilizzo dei contenuti musicali richiedono verifiche
specifiche: la pubblicazione non attribuisce diritti sui brani eseguiti.

La configurazione proposta non registra ordinariamente audio e video.
Qualsiasi registrazione o raccolta educativa aggiuntiva richiede finalità,
procedura e condizioni proprie, distinte dalla semplice partecipazione.
Prima dell'uso si verificano anche i comportamenti effettivi delle
piattaforme esterne: impostare il nodo senza registrazione non dimostra
che un servizio remoto non conservi dati. Inquadrature circoscritte,
stato della connessione visibile e schermata neutra a sessione chiusa
sono requisiti operativi della relazione fra i luoghi.

\section{Tre fasi di sviluppo tecnico, senza scorciatoie geografiche}
\begin{description}
\item[Fase 1: due nodi in condizioni locali controllate.]
Si costruisce la coppia e si verifica l'intera catena: strumenti, audio,
OLED, videocamere, elaborazione e collegamento. Si provano orientamento
quasi speculare, comprensibilità dei gesti, avvio reciproco e arresto.
La prova comprende ripristino e sostituzione di componenti. Il suo risultato
caratterizza quella configurazione; non dimostra ancora un servizio fra città.
\item[Fase 2: collegamento regionale.]
Due luoghi disponibili vengono collegati e misurati nelle fasce previste
per l'uso pubblico. La distanza regionale è un criterio di progressione
organizzativa, non una garanzia di bassa latenza: conta il percorso reale
dei segnali. Le prove musicali stabiliscono quali attività simultanee sono
sostenibili con quella coppia e quali condizioni richiedono sospensione.
\item[Fase 3: lunga distanza da caratterizzare.]
Si studia un ulteriore collegamento dopo le verifiche precedenti,
considerando percorso di rete, sincronizzazione audiovisiva, variabilità,
fusi orari e organizzazione. Non si promette anticipatamente esecuzione
simultanea fluida a qualunque distanza. Le attività alternative si
presentano con il loro nome; non sostituiscono silenziosamente la prova
della prestazione musicale richiesta.
\end{description}

I criteri tecnici del capitolo dedicato all'architettura devono diventare
un verbale per ciascuna coppia: configurazione, luogo, orario, misure,
attività provate, esiti e limiti. Un cambio di connessione, software o
periferica che interessa il percorso critico richiede una verifica mirata.
Per la modalità studio si prova separatamente la disponibilità dei
materiali e il funzionamento della piattaforma scelta, propria o esterna.

\section{Dodici settimane di programma pubblico}
Il pilota proposto comprende dodici settimane, con due appuntamenti
connessi settimanali di due ore. Sono 24 incontri e 48 ore di servizio
congiunto; ciascun nodo ospita le stesse 48 ore. La somma delle ore nei
due luoghi è 96, ma non va presentata come 96 ore distinte di collegamento.
Si aggiunge, per ogni appuntamento e per ogni facilitatore, un'ora di
preparazione e chiusura. Servono quindi 144 ore complessive di facilitazione,
72 per nodo. Trasporti straordinari e formazione iniziale vanno conteggiati
separatamente quando necessari.

Il programma segue le verifiche tecniche e non include automaticamente
le tre fasi geografiche. La ricerca educativa riguarda un sottoinsieme
consenziente e predefinito, con propri criteri e calendario. Incontri,
persone, sessioni di studio e ore di collegamento restano conteggi distinti.

\section{Costi comuni e costi per nodo}
Le seguenti sono \textbf{stime illustrative di pianificazione}, non
preventivi, prezzi verificati sul mercato o finanziamenti disponibili.
Si assumono 55 euro per ora tecnica e 35 euro per ora di facilitazione,
unicamente per rendere riproducibile il calcolo. Imposte, forma contrattuale
e oneri effettivi devono essere definiti nella preventivazione. Le dotazioni
sono accantonamenti per strumenti, OLED, videocamera, audio, calcolo e
involucro: una distinta componenti dovrà sostituirli prima dell'acquisto.

\begin{description}
\item[A --- Coppia per il pilota di dodici settimane.]
Sviluppo comune e documentazione: 240 ore tecniche, 13200 euro.
Dotazione: $2\times6000=12000$ euro. Messa in servizio locale:
$2\times24\times55=2640$ euro. Investimento iniziale: 27840 euro.
L'esercizio comprende 144 ore di facilitazione, 5040 euro;
assistenza comune di 24 ore, 1320 euro; assistenza locale di 12 ore
per nodo, 1320 euro complessivi; logistica e consumabili, 600 euro
per nodo, 1200 euro; servizi condivisi, 600 euro.
Esercizio: 9480 euro. Subtotale: 37320 euro; con riserva del 20\%:
\textbf{44784 euro}.
\item[B --- Coppia per un programma annuale regolare.]
Sviluppo comune: 400 ore tecniche, 22000 euro. Dotazione:
$2\times9000=18000$ euro. Messa in servizio:
$2\times48\times55=5280$ euro. Investimento iniziale: 45280 euro.
Si ipotizzano 36 settimane e due incontri settimanali: 72 incontri,
144 ore di collegamento e 432 ore complessive di facilitazione,
15120 euro. Assistenza comune: 72 ore, 3960 euro; assistenza locale:
36 ore per nodo, 3960 euro complessivi; logistica e consumabili:
1200 euro per nodo, 2400 euro; servizi condivisi: 1200 euro.
Esercizio: 26640 euro. Subtotale del primo anno: 71920 euro;
con riserva del 20\%: \textbf{86304 euro}.
\end{description}

A e B sono scenari alternativi calcolati da zero, non costi cumulativi.
Il passaggio da A a B richiederebbe un computo delle componenti riutilizzabili,
non la somma dei due totali. La fase di lunga distanza non ha qui un prezzo:
richiede un preventivo successivo alla caratterizzazione della connessione.
Una seconda persona di presidio in ciascun nodo raddoppierebbe la sola
voce di facilitazione, non lo sviluppo comune né tutto il budget.

Sono esclusi opere e adeguamenti dei siti, canoni, connettività dedicata,
trasferte straordinarie, ricerca scientifica indipendente, verifiche
specialistiche, permessi, assicurazioni e licenze musicali o didattiche.
Le ore comuni coprono l'integrazione di una configurazione circoscritta;
non stimano una piattaforma educativa proprietaria completa o l'addestramento
di modelli di intelligenza artificiale. Queste voci devono essere
quantificate prima del budget esecutivo. La riserva non le rende incluse.

Il costo totale di possesso deve distinguere parte comune e due nodi:
\[
\mathrm{TCO}_{N}=C_{\mathrm{comune},0}
+\sum_{i=1}^{2} C_{i,0}
+\sum_{t=1}^{N}\left(O_{\mathrm{comune},t}
+\sum_{i=1}^{2}O_{i,t}\right)
+C_{\mathrm{uscita}}-V_{\mathrm{residuo}}.
\]
Gli $O$ comprendono personale, manutenzione, servizi, sostituzioni e oneri
realmente applicabili. Per B, a costi invariati, le sole voci stimate danno
$45280+3\times26640=125200$ euro su tre anni, oppure 150240 euro con
la riserva proposta: un subtotale, non il TCO completo. Il modello
preliminare non applica attualizzazione. Ore volontarie e beni prestati
vanno esposti separatamente: riducono l'esborso, non il lavoro necessario.

\section{Manutenzione, passaggi decisionali e arresto}
Ogni apertura richiede controllo locale e prova fra i nodi. Configurazioni
e ripristino devono essere utilizzabili anche da un secondo operatore.
Gli aggiornamenti sono collaudati sulla coppia; il registro dei guasti
indica nodo, sintomo, intervento ed esito.

\begin{description}
\item[G0: disponibilità della coppia.]
Interlocutori legittimati, condizioni d'uso e risorse iniziali per entrambi
i luoghi. Si cambia sito se un nodo non può rispettare accesso o convivenza.
\item[G1: due nodi controllati.]
Prove della fase locale e criticità di sicurezza risolte prima dell'apertura.
\item[G2: connessione qualificata.]
Prova nel percorso reale, attività musicali validate per quella configurazione,
comandi di interruzione efficaci e responsabilità operative concordate.
\item[G3: programma pubblico.]
Calendario, personale e risorse coperti nei due luoghi. Assenza di presidio,
problemi di sicurezza o perdita delle condizioni necessarie all'attività
sospendono il collegamento; la modalità studio resta disponibile soltanto
se verificata e chiaramente annunciata.
\item[G4: continuità o conclusione.]
Esiti, accessibilità, costi e qualità dell'incontro orientano prosecuzione,
modifica o conclusione. Ritiro dei nodi, riuso e cancellazione dei dati
vanno concordati prima dell'apertura.
\end{description}

